\begin{table}[!htbp]
\caption{Statistics of the held-out scenarios (500 runs per maneuver level).}
\label{tab:S11}
\centering\footnotesize\setlength{\tabcolsep}{5pt}
\begin{tabular}{@{}lccc@{}}
\toprule
Quantity & Low (0.5 g) & Medium (2 g) & High (4 g) \\
\midrule
Peak lateral acceleration, median (g) & 0.49 & 1.95 & 3.91 \\
Peak lateral acceleration, 5th--95th percentile (g) & 0.38--0.61 & 1.54--2.39 & 3.06--4.93 \\
Peak lateral acceleration, median (m/s$^2$) & 4.8 & 19.1 & 38.4 \\
Altitude range (km) & 21.9--56.3 & 22.2--56.2 & 21.6--56.8 \\
Mean number of visible satellites & 5.25 & 5.23 & 5.26 \\
Realized IR outage fraction (\%) & 20.7 & 21.2 & 20.3 \\
Time inside radar coverage (\%) & 40.4 & 39.4 & 37.2 \\
Runs never inside radar coverage (\%) & 0.6 & 1.0 & 3.0 \\
Radar entry time, median (s) & 276 & 284 & 289 \\
Radar entry time, 5th--95th percentile (s) & 184--400 & 193--402 & 187--429 \\
\bottomrule
\end{tabular}

\smallskip
\parbox{\textwidth}{\footnotesize $g = 9.80665$ m/s$^2$. Visible satellites: line of sight above the Earth's surface and range below 4500 km, averaged over IR scan times. Radar entry time: first epoch inside radar coverage, over the runs that enter it.}
\end{table}
